\begin{conf-abstract}
{From residence time to flow rates: insights from \ce{^234 U}/\ce{^238 U} activity ratios and REE patterns of the ancient Kurnub Group aquifer}
{Perach Nuriel\textsuperscript{*1,2}, Daniel Ish-Shalom\textsuperscript{2,3}, Roi Ram\textsuperscript{2}, Christian Siebert\textsuperscript{4}, Boas Lazar\textsuperscript{3}, Ludwik Halicz\textsuperscript{2,5}, Avihu Burg\textsuperscript{2}}
{(1) Department of Earth Sciences, University of Geneva, Genève, CH-1205 Switzerland \\
(2) Geological Survey of Israel, 32 Yeshayahu Leibowitz St., Jerusalem, 9692100 Israel\\
(3) Institute of Earth Sciences, The Hebrew University, Jerusalem, 9190401 Israel\\
(4) Department Catchment Hydrology, Helmholtz-Centre for Environmental Research UFZ, 06120 Halle/Saale Germany\\
(5) Biological \& Chemical Research Centre, Faculty of Chemistry, University of Warsaw, 02-089 Warsaw, Poland}
{Measurements of uranium isotopes and rare earth element (REE) compositions were used to evaluate water-rock interaction processes, residence times, mixing, and flow rates in the Kurnub Group Aquifer (KGA), a deep confined regional sandstone aquifer extending from the Sinai Peninsula (Egypt) to the Negev Desert (Israel). The study involved developing a methodology that allowed to separate both U and REE from the brackish groundwater using a single column. Groundwater was sampled along two main flow paths of the KGA: East (n=9) and North (n=8), and from the overlaying Judea Group Aquifer (JGA; n=2). Overall, samples along the KGA flow paths show consistent low REE concentrations (ppt levels) and patterns, with low LREE and HREE and high MREE. The KGA samples have high 234U/238U activity ratios (AR of up to 8.37), with a systematic exponential decrease in AR downflow as a function of distance and head elevation. This pattern suggests that AR decreases downflow mainly by radioactive decay, enabling the estimation of its long-term average flow rate to 24 cm/yr, similar to recently obtained 81Kr-based estimations. An exceptionally high AR is observed near a major fault zone, suggesting that the highly enriched 234U may originate from this deep-seated active deformation zone, whereby fluid interactions with fractured rocks lead to the preferential leaching of the 234U isotope. Assuming that the computed average flow rate is representative also for the upper parts of the aquifer in Sinai, the uranium isotope-based maximum residence time of groundwater in the KGA is estimated to be ~1.3 Myr.}
\end{conf-abstract}
